% !TEX root = ../monograph.tex
\chapter{Local gauge-invariant observables}
\label{ch02:observables}

A continuum observable needs more than a lattice expression with the right
classical limit. Integration over fluctuations can mix it with other
observables, including the constant. We therefore separate two tasks. This
chapter constructs finite-dimensional spaces of classical polynomial labels
and local lattice representatives of those labels. Later, the
renormalization group will supply the cutoff-dependent linear maps that
normalize the representatives at a fixed physical scale.

The classical connections used below are test configurations. Their jets
provide a precise way to say which local polynomial a lattice observable
represents. We do not assume that a continuum random gauge field is a smooth
connection, or even that such a random connection has been constructed.
Throughout, $G\subset U(n_G)$ and the metric on its Lie algebra $\gLie$ are
those of Section~\ref{ch01:gauge}.

\section{Radial jets and polynomial labels}
\label{ch02:jets-section}

For a smooth connection near the origin, radial gauge is the condition
$x^\mu A_\mu(x)=0$. In this gauge the connection is recovered from curvature
by the usual radial identity
\begin{equation}
 A_\mu(x)=\int_0^1 t x^\nu F_{\nu\mu}(tx)\,dt.
 \label{ch02:radial-identity}
\end{equation}
Repeated coordinate indices are summed from $1$ to $4$. The identity
explains why a homogeneous potential of polynomial degree $k+1$ carries
the same scaling dimension as a $k$th derivative of curvature, namely
$k+2$. We use this observation to define labels algebraically.

\begin{definition}[Radial jets]
\label{ch02:radial-jets}
A formal radial jet is a sequence $J=(A_k)_{k\ge0}$, where
$A_k=(A_{k,\mu})_{\mu=1}^4$ is a $\gLie$-valued one-form with homogeneous
polynomial coefficients of degree $k+1$, satisfying
\[
 \sum_{\mu=1}^4x^\mu A_{k,\mu}(x)=0.
\]
Write $\mathcal J$ for the space of these sequences and
$\mathcal J_{\le d}$ for the finite-dimensional real space obtained by
retaining only $k+2\le d$. It is the zero space for $d<2$.
\end{definition}
Only finitely many jet coordinates enter any polynomial below. A coordinate
of $A_k$ has classical scaling dimension $k+2$ and polynomial degree one.
The scaling dimension of a monomial is the sum of those of its factors.
For a homogeneous polynomial in jet coordinates, we also call its scaling
dimension its \emph{weight}.
The ordinary polynomial degree counts factors of jet coordinates; to
distinguish it from scaling dimension we call it the \emph{factor count}.
Constant gauge transformations act on every coefficient by $\Ad(G)$.
Orthogonal changes of Euclidean coordinates act by the corresponding
one-form transformation. Both actions preserve scaling dimension and factor count.

\begin{definition}[Polynomial labels]
\label{ch02:labels}
For $d\in\NN$, let
\begin{equation}
 V_d=\CC[\mathcal J]^{\Ad(G)}_{\text{scaling dimension}\le d}
 \label{ch02:label-space}
\end{equation}
be the space of gauge-invariant complex polynomials of scaling dimension
at most $d$. Choose nested real bases $(b_{d,i})$ whose elements have definite scaling dimension
$w_i$ and factor count $f_i$, with the constant polynomial as one basis
vector. Conjugation means conjugating coefficients in these real bases.
Define the factor-count operator by
\begin{equation}
 N_{\mathrm{fac}}b_{d,i}=f_i b_{d,i},\qquad
 t^{N_{\mathrm{fac}}}b_{d,i}=t^{f_i}b_{d,i}\quad(t>0).
 \label{ch02:factor-scaling}
\end{equation}
\end{definition}
The spaces are finite-dimensional: only $A_k$ with $k+2\le d$ can occur,
and every nonconstant coordinate has scaling dimension at least two. The invariant
subspace decomposes according to scaling dimension and factor count, so the stated bases can be chosen once and for all. The
constant has bidegree $(0,0)$; the zero-factor subspace is precisely
$\CC\mathbf1$.

The curvature at the origin is
$F_{\mu\nu}=\partial_\mu A_{0,\nu}-\partial_\nu A_{0,\mu}$: the
commutator term vanishes there in radial gauge. Its coordinates have scaling dimension
two and factor count one. An important invariant polynomial is
\begin{equation}
 c_{F^2}=\frac14\sum_{\mu,\nu}
                 \langle F_{\mu\nu},F_{\mu\nu}\rangle_{\gLie},
 \label{ch02:curvature-square}
\end{equation}
with bidegree $(4,2)$. The inner product is extended complex bilinearly when
used in a complex polynomial, not sesquilinearly.

There are three distinct integers in this notation. A basis element has
scaling dimension $w_i$ and factor count $f_i$; it also belongs to a chosen space
$V_d$ with $w_i\le d$. We shall retain that last choice as part of the
observable label, writing $v=(d,c)$ with $c\in V_d$. The same polynomial
$c$ can occur at different tags $d$. Its eventual normalization depends on
the tag, so the resulting insertions must not be identified without a
comparison theorem. None of these integers is being asserted to be the
exact scaling dimension of the interacting quantum field.

\section{Test holonomies and Taylor normalization}
\label{ch02:taylor-section}

A label is a polynomial of a jet. A lattice observable is a function of
link matrices. To compare them, insert the parallel transports of a
truncated radial connection into the lattice observable.

For $J\in\mathcal J_{\le d}$, write
$A^{\le d}_\mu(x)=\sum_{k+2\le d}A_{k,\mu}(x)$. For a positive link
$(y,\mu)$ of $\ZZ^4$ and a real parameter $h$, define
\begin{equation}
 U(h,J)(y,\mu)=u(1),\qquad
 u'(t)=u(t)h A^{\le d}_\mu(h(y+te_\mu)),\quad u(0)=1.
 \label{ch02:test-holonomy}
\end{equation}
The solution remains in $G$. Standard dependence of solutions of analytic
ordinary differential equations on parameters gives analytic dependence on
the finitely many jet coordinates near zero; see
\cite{Teschl2012}. Reversed links have inverse transport, in the convention
of Section~\ref{ch01:wilson-section}.

Define the jet dilation
\begin{equation}
 \operatorname{sc}_hJ=(h^{k+2}A_k)_k.
 \label{ch02:jet-dilation}
\end{equation}
Homogeneity in \eqref{ch02:test-holonomy} gives
$U(h,J)=U(1,\operatorname{sc}_hJ)$. In particular, a homogeneous label of
scaling dimension $w$ satisfies $c(\operatorname{sc}_hJ)=h^wc(J)$.

\begin{definition}[Taylor map and local section]
\label{ch02:taylor-map}
For a finite-link function $O$ analytic near the identity configuration,
$\operatorname{Tay}_d O$ is the part of scaling dimension at most $d$ in
the Taylor expansion at $J=0$ of $O(U(1,J))$.
A local polynomial section is a complex-linear map
\[
 P_d^{\mathrm{loc}}:V_d\longrightarrow
       \{\text{polynomial functions of finitely many link entries and their
       conjugates}\}
\]
with $\operatorname{Tay}_dP_d^{\mathrm{loc}}[c]=c$.
\end{definition}
The term section simply means a right inverse of the Taylor map. The
polynomial condition is global in the link matrices; it is stronger than
analyticity near the identity. For a gauge-invariant observable the Taylor
polynomial is $\Ad(G)$-invariant, since constant gauge transformations
conjugate the test holonomies.

Equivalently, the section identity requires, for every $J$ and $0\le q\le d$,
\begin{equation}
 \left.\frac{d^q}{dh^q}P_d^{\mathrm{loc}}[c](U(h,J))\right|_{h=0}
 =\left.\frac{d^q}{dh^q}c(\operatorname{sc}_hJ)\right|_{h=0}.
 \label{ch02:taylor-identity}
\end{equation}
It fixes every homogeneous component of scaling dimension at most $d$. It also does not prescribe the value of the lattice observable
on a typical rough configuration.

\section{Finite-stencil representatives}
\label{ch02:sections}

The construction uses lattice curvature transported to a common point.
Extend the metric \eqref{ch01:metric} to the ambient real matrix space as
$\operatorname{Re}\tr(X^*Y)/n_G$, and let $\Pi_{\gLie}$ be orthogonal
projection onto $\gLie$. For a plaquette based at $z$, set
\[
 F_\square(U)=\Pi_{\gLie}\frac{U_\square-U_\square^{-1}}2.
\]
If $\gamma$ is a lattice path from the origin to $z$, its transported
curvature is
\begin{equation}
 \sigma_{\gamma,\square}(U)=\Ad(U_\gamma)F_\square(U).
 \label{ch02:transported-curvature}
\end{equation}
Under a site gauge transformation it is conjugated by the gauge element
at the origin. It is a real polynomial in the entries and conjugate
entries of the links it reads, because inverse unitary matrices are
conjugate transposes.

\begin{theorem}[Local polynomial representatives]
\label{ch02:local-sections}
For each $d$ there is a finite centered stencil $\Sigma_d\subset\ZZ^4$,
invariant under the signed permutation group $H_4$, and a local polynomial
section $P_d^{\mathrm{loc}}$ with the following properties.
Its values are gauge invariant and local on $\Sigma_d$. It preserves the
unit and complex conjugation, and satisfies
\begin{equation}
 P_d^{\mathrm{loc}}[c](h^*U)=P_d^{\mathrm{loc}}[h\cdot c](U),
 \qquad h\in H_4.
 \label{ch02:section-covariance}
\end{equation}
Here $h^*$ pulls back lattice links and
$(h\cdot c)(J)=c(h^{-1}\cdot J)$ is the induced action on labels.
Moreover, each basis column $P_d^{\mathrm{loc}}[b_{d,i}]$ is a polynomial
in finitely many transported curvatures, and every monomial in that
polynomial has at least $f_i$ factors.
\end{theorem}
The last assertion is what makes the section useful for estimates. An
arbitrary right inverse of the Taylor map need not preserve this lower
bound on factor count.

\begin{proof}
We give the finite-dimensional construction, invoking the usual inverse
function theorem rather than reproving it.

First consider the maps $J\mapsto\sigma_{\gamma,\square}(U(1,J))$.
Their linearizations at the zero jet separate $\mathcal J_{\le d}$.
Indeed, differentiating transport at the identity gives the line integral
of the connection along the link. The derivative of a plaquette curvature
is therefore its circulation. The derivative of the conjugating transport
in \eqref{ch02:transported-curvature} contributes nothing, since the
curvature itself vanishes at the identity.

To verify separation, suppose all these linearized plaquette measurements
vanish on $A=\sum_{k=0}^{d-2}A_k$. For each coordinate plane the circulation
around its unit square based at $z$ is the polynomial
\[
 N_{\mu\nu}(z)=\int_0^1\!\int_0^1
 (\partial_\mu A_\nu-\partial_\nu A_\mu)
                  (z+te_\mu+ue_\nu)\,dt\,du.
\]
It vanishes at every integer $z$, and hence is the zero polynomial. If $k$
is the largest degree index with $A_k\ne0$, its homogeneous part of degree
$k$ is $\partial_\mu A_{k,\nu}-\partial_\nu A_{k,\mu}$. Thus $dA_k=0$.
The radial identity and Euler's identity give
\[
 \sum_\mu x^\mu
  (\partial_\mu A_{k,\nu}-\partial_\nu A_{k,\mu})=(k+2)A_{k,\nu}.
\]
It follows that $A_k=0$, a contradiction. Descending through the degree
indices proves separation. Since $\mathcal J_{\le d}$ is finite-dimensional,
a finite collection of measurements already has injective joint derivative.

Fix such a collection and write
\[
 \mathscr S(J)=\bigl(\sigma_1(U(1,J)),\ldots,
                          \sigma_M(U(1,J))\bigr)\in\gLie^M,
 \qquad D=D\mathscr S(0).
\]
The map is analytic, $\mathscr S(0)=0$, and $D$ is injective. Both the jet
space and $\gLie^M$ have invariant real inner products. Hence the inverse
of $D$ on its range, composed with orthogonal projection onto that range,
is an equivariant linear map $C$ with $CD=1$. The inverse function theorem
applied to $C\mathscr S$ gives an analytic local inverse $g$; this is the
standard finite-dimensional theorem, for example
\cite[Theorem~9.24]{Rudin1976} together with its analytic form.
Set $f=gC$. Then $f(0)=0$ and $f\mathscr S(J)=J$ near zero.
Equivariance follows from uniqueness of the inverse as a germ.

Let $T_df$ be the ordinary Taylor polynomial of $f$ through degree $d$.
It has no constant term and is equivariant. For each real basis polynomial
$b_{d,i}$ define $Q_i(z)=b_{d,i}(T_df(z))$. Then $Q_i$ is invariant, is real
on real arguments, and has no monomial of degree below $f_i$. The unit
column is one. Consequently
\[
 \Psi_c(U)=\sum_i c_iQ_i(\sigma_1(U),\ldots,\sigma_M(U))
 \quad\text{for }c=\sum_i c_i b_{d,i}
\]
is a gauge-invariant polynomial cylinder function. Since
$T_df(\mathscr S(J))=J+O(|J|^{d+1})$, its Taylor map is $c$.
Each jet coordinate has scaling dimension at least two, so every monomial
of polynomial degree greater than $d$ also has scaling dimension greater
than $d$. No truncation of $Q_i$ after substitution is
needed.

Finally average over the finite lattice symmetry group:
\begin{equation}
 P_d^{\mathrm{loc}}[c](U)
   =\frac1{|H_4|}\sum_{h\in H_4}\Psi_{h^{-1}\cdot c}(h^*U).
 \label{ch02:section-average}
\end{equation}
Uniqueness of the parallel-transport equation gives
$h^*U(h_0,J)=U(h_0,h^{-1}\cdot J)$ for a scale parameter $h_0$.
Thus every summand still has Taylor map $c$. Reindexing the finite sum
proves \eqref{ch02:section-covariance}. Reality and the unit are preserved.
A signed image of a transported curvature is another such curvature,
possibly with a minus sign from orientation reversal. The group action
preserves factor count, so the lower bound on monomial degree survives
averaging. Enclosing all paths and plaquettes in a centered box produces
the finite stencil $\Sigma_d$.
For $d<2$ there are only constants and the construction reduces to the
constant section.
\end{proof}

All choices in this theorem are made on the covering lattice $\ZZ^4$,
before selecting a torus period. For a site $y\in\Lambda_s$, write
$P_d^{\mathrm{loc}}[c](y;U)$ for evaluation after translating the stencil
to $y$ and reducing its links modulo $s$. The construction remains
well-defined if different covering links become the same torus link.
Gauge invariance, covariance and the polynomial identities survive that
identification.

The abstract jet space admits an $O(4)$ action, but the covariance in
\eqref{ch02:section-covariance} concerns only $H_4$. Arbitrary rotations do
not act on the lattice links. This distinction prevents the classical
label symmetry from being mistaken for an already proved symmetry of the
continuum correlations.

\section{Small-field bounds for the representatives}
\label{ch02:bounds}

Recall that $V_d$ is the polynomial label space of scaling dimension at most $d$ from Definition~\ref{ch02:labels}; factor count counts curvature-jet factors instead of their scaling dimensions.

We now use the factor-count property of
Theorem~\ref{ch02:local-sections}. Recall from \eqref{ch01:chord} that
$\chi(u)=\|u-1\|_{\mathrm F}/\sqrt{n_G}$ measures the size of a plaquette.
On a fixed local stencil, a bound $\chi(U_\square)\le t$ gives a factor
$t$ for each transported curvature in \eqref{ch02:transported-curvature}.
The same factor is retained under suitably scaled derivatives.

For a collection $A=(A_e)_{e\in E_s}$ of Lie algebra elements, define the
left derivative
\[
 D_A^L F(U)=\left.\frac{d}{d\varepsilon}
       F\bigl((\exp(\varepsilon A_e)U_e)_{e\in E_s}\bigr)
          \right|_{\varepsilon=0}.
\]
An ordered list of these derivatives is used; no commutation of Lie
derivatives is assumed. Fix a nonnegative spatial exponent $p_0$ and a
site $y$. A permitted direction at threshold $t$ satisfies
\begin{equation}
 |A_e|_{\gLie}\le t\bigl(1+\operatorname{dist}_s(e^-,y)\bigr)^{p_0},
 \label{ch02:direction}
\end{equation}
where $e^-$ is the initial site of the positive link. This is a local
definition for the estimate; the later polymer norms will specify which
orders and spatial exponents they retain.

\begin{proposition}[Uniform section bounds]
\label{ch02:section-bounds}
For every $d,q,p_0$ there are finite constants $C_{d,q,p_0}$ and $C_d$,
independent of $s$, $y$ and $0<t\le1$, with the following properties.
If every plaquette used by the translated section has distance from the identity at most $t$,
then for every list of permitted directions of length $k\le q$,
\begin{equation}
 \left|D_{A_1}^L\cdots D_{A_k}^L
       P_d^{\mathrm{loc}}[b_{d,i}](y;U)\right|
       \le C_{d,q,p_0}t^{f_i}.
 \label{ch02:section-small-bound}
\end{equation}
For arbitrary configurations,
\begin{equation}
 \sup_{i,y,U}|P_d^{\mathrm{loc}}[b_{d,i}](y;U)|\le C_d.
 \label{ch02:section-global-bound}
\end{equation}
The same constants apply when the stencil wraps around the torus.
\end{proposition}
\begin{proof}
Projection is contractive and $U^{-1}=U^*$, so
\[
 |F_\square(U)|_{\gLie}
 \le\frac{\|U_\square-U_\square^*\|_{\mathrm F}}{2\sqrt{n_G}}
 \le\chi(U_\square).
\]
Transport by a unitary matrix preserves the metric. An undifferentiated
curvature measurement is therefore bounded by $t$ on the stated domain.

On the fixed stencil the right side of \eqref{ch02:direction} is bounded
by a constant times $t$, uniformly in the period: periodic distance is
no larger than distance to a chosen covering-space representative.
Differentiating a fixed path product $k$ times gives a finite sum of
products with exactly $k$ direction matrices and otherwise unitary
factors. Its derivatives are bounded by a constant times $t^k$.
The same conclusion holds for each transported curvature. For a positive
order $k$, this is in particular bounded by a constant times $t$, since
$t\le1$.

Every monomial representing the column $b_{d,i}$ has at least $f_i$
curvature factors. Apply the ordered product rule: each factor, whether
or not differentiated, contributes a factor bounded by a constant times
$t$. Summing the finitely many resulting terms proves
\eqref{ch02:section-small-bound}. Matrix entries of a unitary matrix have
modulus at most one, so the polynomial cylinder formula gives
\eqref{ch02:section-global-bound}.
If covering links coincide on the torus, one first differentiates the
covering expression and assigns the same direction to identified links.
The number of terms and their bounds do not increase with the period.
\end{proof}

The scaling operator $t^{N_{\mathrm{fac}}}$ from
\eqref{ch02:factor-scaling} is tailored to this estimate. By linearity,
permitted derivatives of $P_d^{\mathrm{loc}}[t^{-N_{\mathrm{fac}}}c]$
are bounded by $C_{d,q,p_0}\sum_i|c_i|$ on the small-field domain. The
powers $t^{f_i}$ cancel column by column. Outside that domain the global
bound instead gives $C_d\sum_i|c_i|t^{-f_i}$. The missing small-field gain
must later be supplied by the probability cost of a large-field region;
it cannot be obtained by extrapolating \eqref{ch02:section-small-bound}.

\section{Local representatives of the equations of motion}
\label{ch02:equation-section}

The ordinary section $P_d^{\rm loc}$ reproduces Taylor coefficients,
but need not reproduce equations of motion. We construct a second
section for this purpose. Everything in this section concerns the
finite-volume Wilson measure; no renormalization estimate is used.

A radial polynomial potential $A_j$, homogeneous of degree $j+1$,
is a \emph{Maxwell solution potential} if
\[
 \sum_\mu\partial_\mu(dA_j)_{\mu\nu}=0\quad\hbox{for every }\nu.
\]
Its curvature is automatically closed. Let $\mathcal I_d\subset V_d$
be the polynomials vanishing on all such solution jets through degree
$d$. This is the degree-$d$ part of the equation ideal; multiplication
is understood only when the resulting scaling dimension remains at most $d$.
The condition preserves both scaling dimension and factor count.
It is the kernel of restriction to solution jets, rather than a
statement about the interacting measure.

\begin{theorem}[An equation-adapted local section]
\label{ch02:equation-adapted}
There are a real, symmetry-equivariant linear section
$\mathcal V_d$ from $V_d$ to local continuous gauge-invariant functions of links, a real equivariant automorphism
$R_d$ of $V_d$, and finite constants $r_d,M_d$, with the following
properties:
\begin{enumerate}
\item $\mathcal V_d1=1$, $R_d1=1$, and
$\operatorname{Tay}_d\mathcal V_d p=p$.
The section has a fixed finite stencil, independent of the torus.
\item If $p$ has factor count $s$ and scaling dimension $w$,
then $R_dp-p$ is a sum of terms whose pair
$(\text{factor count},\text{scaling dimension})$ is strictly larger than $(s,w)$
in lexicographic order. Every term of $\mathcal V_dp$ is a product
of at least $s$ transported curvature coordinates.
\item For every $p\in\mathcal I_d$, every torus of side at least $M_d$,
every anchor $o$, every real $\beta\ne0$, and every bounded measurable
$S$ depending only on links with both endpoints farther than $r_d$
from $o$,
\begin{equation}
 \int \mathcal V_d(R_dp)(o,U)S(U)\,d\mu_\beta(U)=0.
 \label{ch02:microscopic-null}
\end{equation}
No gauge invariance or regularity of $S$ is assumed.
\end{enumerate}
The small-field derivative bounds for a fixed section proved above
also apply to $\mathcal V_d$, with its own finite constants.
\end{theorem}

The regrading $R_d$ is essential: the exact interacting relation is
$\mathcal V_d(R_d\mathcal I_d)$, not an assertion that
$P_d^{\rm loc}(\mathcal I_d)$ has zero Wilson correlations.
The proof below supplies the finite-stencil construction, whereas the
compact-group integration by parts underlying it is the usual lattice
equation-of-motion identity for Wilson's measure
\cite{Wilson1974,OsterwalderSeiler1978}.

\subsection{Curvature tests avoiding a prescribed link}

For $u\in\gLie$, pairing the transported curvature
\eqref{ch02:transported-curvature} with $u$ gives the real-valued function
$U\mapsto\langle u,\sigma_{\gamma,\square}(U)\rangle_{\gLie}$.
We call any such function a curvature coordinate.
Call it admissible for an oriented link $e$ if neither its plaquette
nor its transporting path uses the underlying unoriented link $e$.
Such a coordinate is independent of $U_e$.

\begin{lemma}\label{ch02:avoiding-tests}
For each $e$, the linearizations at the zero jet of admissible
curvature coordinates span the dual of the finite radial jet space
$\mathcal J_{\le d}$.
\end{lemma}
\begin{proof}
It suffices to show that their common annihilator is zero.
Removing one link does not disconnect $\ZZ^4$, so every plaquette
not containing $e$ admits a transporting path avoiding $e$.
Let $A$ be the polynomial one-form representing a jet in the annihilator.
At the zero connection the derivative of transport around an oriented
unit plaquette $\square$ in direction $A$ is
$\oint_{\partial\square}A=\int_\square dA$ by Stokes' theorem.
Transport to the anchor does not change this first variation.
Consequently $\int_\square dA=0$ for every plaquette that does not
contain $e$.

For each orientation $(\mu,\nu)$, this integral is a polynomial in
the base point $z$ of the plaquette. It vanishes at all but finitely
many integer points and hence identically, by the one-variable polynomial
identity applied successively in each coordinate. Let $A_j$ be the
highest nonzero homogeneous layer of $A$. Averaging a polynomial over
a translated unit square preserves its highest homogeneous part as a
polynomial in $z$. Thus the degree-$j$ part of the integral is
$(dA_j)_{\mu\nu}(z)$. Since the integral vanishes for every orientation,
$dA_j=0$. Radial gauge and Euler's identity give
\[
 0=\partial_i(x^\mu A_{j,\mu})
   =A_{j,i}+x^\mu\partial_i A_{j,\mu}
   =(j+2)A_{j,i},
\]
where closedness was used in the last equality.
This contradicts the choice of $j$. Descending through the layers
proves the assertion. Finite-dimensional duality now selects a
finite spanning collection of admissible coordinates.
\end{proof}

\subsection{Construction of the equation-adapted section}

\begin{proof}[Proof of Theorem~\ref{ch02:equation-adapted}]
We first represent the generators of the equation ideal by Wilson forces.
We then integrate by parts in a single link and assemble the local section.

\smallskip\noindent\emph{Step 1. Representation by Wilson forces.}

Write $S_W$ for the Wilson action without its factor $\beta$.
For a positive link $e$ and a direction $X\in\gLie$,
the force $F_e\in\gLie$ is characterized by
\[
 \partial_{e,X}S_W=\langle X,F_e\rangle,\qquad
 \partial_{e,X}f(U)=
 \left.\frac{d}{dt}f(U_e\leftarrow e^{tX}U_e)\right|_{t=0}.
\]
Differentiating the plaquettes incident on $e$ expresses $F_e$
as a signed sum of plaquette-curvature coordinates based at the
initial vertex of $e$. The precise signs follow the chosen
plaquette orientations. Transport the force to the anchor by a
path avoiding $e$, and pair it there with an arbitrary element of $\gLie$.
At the zero connection the force is the signed sum, to first order in
$A$, of the integrals $\int_\square dA$ over plaquettes incident on $e$.

Fix a layer $j$. These linear force functionals, restricted to jets
through layer $j$, span every functional which acts only on $A_j$
and annihilates the Maxwell solution potentials in that layer.
Indeed, if all force functionals vanish, their values at integer
base points form identically zero polynomials. Their top-degree
terms are
\[
 -\sum_\mu\partial_\mu(dA_j)_{\mu\nu}.
\]
For a link in direction $\nu$, the two incident $(\mu,\nu)$ plaquettes
enter with opposite signs. Their integrals therefore give a first
difference in the $\mu$ direction. Thus the top layer is a Maxwell
solution potential. Every functional of the indicated type
annihilates the common force kernel and is therefore a finite
linear combination of these forces. Notice that this argument
identifies the top layer only; it does not identify the entire
kernel of the discrete force system with the continuum solution
jet space.

Here is the corresponding description of the polynomial ideal.
Choose a linear projection of each jet layer onto its solution
subspace and let $\mathcal E$ be their direct sum. In the full,
not necessarily invariant, polynomial algebra the kernel of
$p\mapsto p(\mathcal EJ)$ is generated by linear functionals
vanishing on those solution subspaces. To check this without
an ideal-generation theorem, use
\[
 \prod_{\nu=1}^s u_\nu-\prod_{\nu=1}^s(\mathcal E^*u_\nu)
 =\sum_{i=1}^s
 (u_i-\mathcal E^*u_i)
 \prod_{\nu<i}\mathcal E^*u_\nu
 \prod_{\nu>i}u_\nu .
\]
Choose homogeneous linear coordinates, so both gradings are
preserved. Every homogeneous element of $\mathcal I_d$ is
therefore a sum of products $\lambda z_1\cdots z_{s-1}$,
where $\lambda$ is a one-layer equation functional.
Individual summands need not be gauge invariant.

Represent $\lambda$ by the finite force combination just constructed
on its truncated jet space. For each force summand involving $e$,
represent each $z_i$ by a linear combination of curvature coordinates
avoiding $e$, using Lemma~\ref{ch02:avoiding-tests} on the \emph{full}
jet space through degree $d$. Multiply these representatives by
that force summand. The resulting function $f$ has at least $s$
curvature factors. Its Taylor polynomial is
\begin{equation}
 \operatorname{Tay}_d f
   =\lambda z_1\cdots z_{s-1}
       +\{\hbox{terms strictly higher in }(s,w)\}.
 \label{ch02:equation-triangle}
\end{equation}
There are two distinct errors. The nonlinear part of a curvature
or force coordinate increases factor count. The linear force
representation was fixed only through the layer of $\lambda$;
its remaining linear terms increase the scaling dimension while preserving the factor
count. This is why the order in
\eqref{ch02:equation-triangle} puts factor count first.

\smallskip\noindent\emph{Step 2. Integration by parts and construction of the section.}

Every force term above has the form
$H\langle X,F_e\rangle$, where $X$ is the fixed direction at the
anchor transported back to the link, and $H$ is the product of the
remaining curvature coordinates and the exterior spectator.
Freeze the other links. Both $X$ and $H$ are then constant in $U_e$. Haar invariance
on the compact group gives
\[
 0=\int_G\partial_{e,X}\{e^{-\beta S_W}\}\,dU_e
   =-\beta\int_G\langle X,F_e\rangle
                          e^{-\beta S_W}\,dU_e.
\]
Multiplication by $H$ is legitimate because it is constant in
$U_e$. Fubini then proves the desired zero expectation, even for a
merely bounded measurable exterior spectator. This proof never
differentiates that spectator.

All lists used in constructing a basis are finite. Enlarge one
radius to contain their plaquettes, paths, and differentiated links.
Enlarge one period threshold so that distinct relevant cover links
cannot become identified with a differentiated link after
periodization. These two choices work for every anchor by translation.
They are independent of $\beta$ and of the exterior spectator.

Average the constructed functions over the gauge element at the
anchor. Every transported curvature transforms there by the adjoint
action, so this is a finite-dimensional average of tensors under the adjoint action.
It preserves the number of curvature factors and
\eqref{ch02:equation-triangle}. It also preserves the exact zero
identity: gauge invariance of the Wilson law transfers the gauge
transformation to the bounded exterior spectator without changing
its support. The result is fully gauge invariant.

Choose a basis of each homogeneous part of $\mathcal I_d$, perform
this construction there, and extend to a basis of $V_d$ by using
$P_d^{\rm loc}$ on a complement. The resulting linear map
$\mathcal F_d$ preserves the unit. Its Taylor map is
\[
 R_d=\operatorname{Tay}_d\mathcal F_d=I+T_d,
\]
where $T_d$ strictly increases the finite lexicographic grading.
Thus $T_d$ is nilpotent and
$R_d^{-1}=I-T_d+T_d^2-\cdots$ is a finite sum.
Average $\mathcal F_d$ simultaneously on domain and range over
signed coordinate permutations and over conjugation. These
operations preserve the ideal, both gradings, and the unit;
hence the same triangular conclusion remains valid.
Finally put
\[
 \mathcal V_d=\mathcal F_dR_d^{-1}.
\]
This gives the Taylor identity, all symmetries, and
$\mathcal V_dR_d p=\mathcal F_dp$ for $p\in\mathcal I_d$.
Both $R_d$ and its inverse preserve the lower bound on factor
count. The fixed finite products therefore obey exactly the
small-field derivative argument used for $P_d^{\rm loc}$,
with new constants for the new stencil.
\end{proof}

The identity \eqref{ch02:microscopic-null} is exact at each cutoff.
It gives no bound on a renormalized coefficient change.
To use it for continuum null fields one must compare this section
with the section used along the actual renormalization history,
with one invertible comparison map for all labels and all spectators.


\section{Examples and the remaining normalization problem}
\label{ch02:normalization}

The unit illustrates the simplest section: $P_d^{\mathrm{loc}}[\mathbf1]=1$
for every $d$. Its eventual physical normalization still depends on the
normalization index. There will be a reference scale for each $d$, and a choice of
coefficient making the renormalized insertion exactly one. This does not
follow by putting a universal factor $a_N^{-d}$ in front of the raw
constant section.

For the curvature-square polynomial $c_{F^2}$ in
\eqref{ch02:curvature-square}, a useful classical representative is the
clover average
\begin{equation}
 P^{\mathrm{cl}}(y;U)=\frac14\sum_{\square\ni y}w(U_\square),
 \label{ch02:clover}
\end{equation}
where the sum has four corner plaquettes in each of the six coordinate
planes, hence $24$ terms. Orientation does not matter because $w(u^{-1})=w(u)$.
For a smooth test connection the expansion
\eqref{ch01:classical-curvature} yields
\[
 P^{\mathrm{cl}}(y;U(h,J))
     =h^4 c_{F^2}(J)+O(h^5)
\]
at the origin, and thus $\operatorname{Tay}_4P^{\mathrm{cl}}=c_{F^2}$.
The factors are consistent: each plaquette contributes
$\frac12h^4|F_{\mu\nu}|_{\gLie}^2$, four copies occur per plane, and the
prefactor is $1/4$. This example fixes the normalization, but the general
construction is free to use a different section satisfying the same
Taylor identity.

The clover formula does not solve the quantum normalization problem.
Under blocking, the curvature square can mix with the constant. A general
renormalized version therefore involves both a multiplicative coefficient
and an additive constant, with coefficients depending on the cutoff and
volume. More generally, a label in $V_d$ can mix with other labels in the
same finite-dimensional space. The correct normalization is consequently
an invertible linear map on $V_d$, whose existence and estimates must be
proved along the chosen renormalization trajectory.

There is also a distinction between polynomial relations and vanishing
fields. Some polynomial combinations represent free Maxwell equations of
motion. The interacting lattice equations arise instead from Haar
integration by parts, and include nonlinear corrections. To identify the
null space of continuum observables requires both those identities and a
rank estimate. Neither the section theorem nor its classical Taylor
normalization identifies the interacting null space with the unmodified
free equation ideal.

The output of this chapter is therefore concrete but limited: a common
finite-stencil representative for every polynomial label, with symmetry
and threshold estimates uniform in the torus period. These are exactly
the local data needed to begin the renormalization of observables. The
uniform control of their correlations and their continuum normalization
remain the task of the subsequent construction.
